% ---------------------------------------------------------------------------------------------
% Rebuilt replacement for the subsection currently titled
%   "Localized temporal processing achieves the most balanced fidelity across paradigms"
% (template/result.tex, line 151 onward).
%
% Why it needed rebuilding, in short:
%   1. The subsection's central quantitative claim, that PatchTST scores $1.00$ on the Markov
%      dimension, is reversed by the redesigned paradigm. PatchTST now ranks last of fifteen
%      models on switching fidelity at two of the three dwell times.
%   2. MICN's frontier status was inflated by decoder-input leakage in the published runs.
%   3. PatchTST is architecturally hybrid, so it is weak evidence for a claim about localized
%      processing specifically.
%   4. The real-data results give a direct test of the claim, which the original text predates.
%
% Numbers below are traceable to analysis/p4_markov/markov_metrics.csv,
% analysis/p1_real/supp_table_real.csv and analysis/p2_baselines/ACCEPTANCE_FAILURE.md.
%
% PENDING: the Pareto frontier itself cannot be regenerated until the noise and shift sweeps are
% re-run with the full sixteen-model roster. Paragraph 3 is written so that it states only what
% the completed paradigms support; the bracketed sentence is the single place where the
% regenerated frontier numbers must be inserted.
% ---------------------------------------------------------------------------------------------

\subsection*{No architecture maintains fidelity across all paradigms, and the advantage of localized processing is paradigm-dependent}

The preceding analyses each isolate a single perturbation dimension. In practice, a forecasting
model deployed for patient monitoring faces noise, frequency shifts, and state changes
simultaneously, and fidelity varies not only across individual dynamical properties but also
across their interactions under combined perturbation \cite{clifford2006advanced,
sornmo2005bioelectrical}. Strong performance on any single paradigm does not guarantee overall
suitability. We therefore ask whether any architecture maintains fidelity across all five
paradigms, and whether the localized-processing hypothesis survives that test.

It does not survive in the form we previously stated. Under the redesigned switching paradigm,
in which the expected dwell time is controlled explicitly rather than the per-sample transition
probability, PatchTST ranks last of the fifteen models on switching fidelity at expected dwell
times of $2$ s and $5$ s, and thirteenth at $10$ s. Its predicted mean dwell time is $29.6$ s
against a true $3.1$ s at $D = 2$ s, and $292.8$ s against a true $11.2$ s at $D = 10$ s: its
forecasts almost never change state. The seasonal-naive floor, which performs no learning,
attains the lowest transition-matrix divergence of any model at $D = 10$ s
($0.004$ nats per step, against $0.108$ for PatchTST). Our earlier conclusion that PatchTST
recovers switching dynamics was an artifact of the previous parameterization, in which the chain
switched with probability $p$ at every sample and the mean dwell time was shorter than one
carrier cycle, so that the decoded state sequence carried little information about switching.

The ordering by pointwise error and the ordering by switching fidelity are close to opposite.
The models with the lowest mean absolute error on this paradigm, MICN (mean), MICN (regre) and
TSMixer, occupy ranks eight to eleven of fifteen on divergence, while the models that reproduce
the switching statistics best, the linear family and the naive floor, rank ninth to fifteenth on
error. A deterministic predictor minimizes expected pointwise error on a stochastic switching
task by declining to commit to a state, which produces a smooth forecast with few transitions,
and only a metric defined on the switching process makes that visible.
[Insert the regenerated Pareto frontier here once the noise and shift sweeps have been re-run
with the full sixteen-model roster; the frontier reported previously rests on MICN scores that
were inflated by decoder-input leakage, as described in Methods.]

The real-data experiments provide a direct test of the localized-processing hypothesis that the
synthetic paradigms alone cannot give. Among trained models the ordering by fidelity does favor
local receptive fields on photoplethysmography and electrocardiography, where ModernTCN, the two
MICN variants and PatchTST lead. It does not hold on electroencephalography: on Sleep-EDF the
decomposition and linear families lead, and the separation between the best local model and the
best linear model is $0.004$ in normalized mean absolute error, which is within the spread across
seeds. The hypothesis is therefore conditional on the signal having a dominant local period, not
general.

Two further observations constrain the claim. First, PatchTST is architecturally hybrid: its
input stage is local, since the series is cut into patches that are projected independently, but
its predictive stage applies full self-attention across all patches. Evidence drawn from PatchTST
alone cannot separate the contribution of local tokenization from that of global mixing, and we
report the bias-group analysis with PatchTST assigned to each group in turn in the Supplement.
Second, on clean photoplethysmography a seasonal-naive forecaster that repeats the last dominant
period outperforms all fifteen trained architectures on every metric we computed, including mean
absolute error ($0.518$ against $0.660$ for the best trained model), phase error
($33.8^{\circ}$ against $56.1^{\circ}$) and heart-rate error ($4.3$ against $4.8$ beats per
minute). Any claim that a particular inductive bias confers fidelity on real periodic signals
must first account for the fact that periodic extrapolation from a five-second history is not
beaten by any of them.

Taken together, these results do not support a recommendation of localized temporal processing as
a general design principle for physiological forecasting. They support three narrower statements.
Local receptive fields confer an advantage where the signal has a dominant local period, and that
advantage disappears on electroencephalography. No architecture in this benchmark, local or
otherwise, reproduces stochastic switching statistics, and the architectures that score best by
pointwise error are among the worst on this dimension. And on the clean periodic signals where
the locality advantage is clearest, none of the trained architectures improves on a
parameter-free periodic baseline, which sets the floor any future architectural claim must clear.
